\documentclass[12pt]{article}
\usepackage[utf8]{inputenc}
\usepackage[T1]{fontenc}
\usepackage{hyperref}
\title{Kalou Language Resources for K-12 Curriculum Development}
\author{Compiled by Claude Sonnet 5.5 (claude-sonnet-5-5) for Andr\'as Kornai}
\date{}
\begin{document}
\maketitle

\section{Introduction}
Kalou (also called Yawa; ISO 639-3 ywa) is a Sepik-family language of Papua New Guinea, closely related to Amal. Speaker counts vary by source: about 1,400 (Wikipedia, 2003), 820 (Verbix), and 2,600 people (Joshua Project). Joshua Project describes the language as unwritten. Searches found no educational materials and no dedicated linguistic description.

\section{Primary Educational Resources}
None found.

\section{Linguistic and Reference Materials}
\begin{itemize}
\item Wikipedia entry on Kalou: \url{https://en.wikipedia.org/wiki/Kalou_language}
\item Wikipedia entry on the related Amal language: \url{https://en.wikipedia.org/wiki/Amal_language}
\item Endangered Languages Project entry: \url{https://www.endangeredlanguages.com/lang/10631}
\item Verbix endangered-languages entry: \url{https://docs.verbix.com/EndangeredLanguages/Kalou}
\end{itemize}

\section{Religious and Community Resources}
Joshua Project profile for the Kalou people. It notes that some audio recordings exist and that some speakers also use Tok Pisin: \url{https://joshuaproject.net/people_groups/12455}

\section{Digital and Audio Resources}
Joshua Project mentions audio recordings but the search results gave no direct links to them.

\section{Curriculum Development Context}
The language is unwritten and not known to be taught in schools. A curriculum would need new community-led documentation, with an orthography and basic descriptive work, before any teaching materials could be made.

\end{document}